\section{Technische und methodische Grundlagen}
\label{sec:grundlagen}

Dieses Kapitel führt nur die Verfahren ein, die für Architektur,
Versuchsdesign und Interpretation der Pipeline notwendig sind. Die Arbeit
entwickelt keine neuen Segmentierungs-, SfM- oder Splatting-Algorithmen. Ihre
wissenschaftliche Leistung liegt in der konsistenten Verbindung und Prüfung
dieser Verfahren.

\subsection{Promptbare temporale Segmentierung}

Das Segment Anything Model überführt visuelle beziehungsweise sprachliche
Prompts in Objektmasken \parencite{kirillov2023segment,metaSam3}. Für die
Pipeline ist die Promptbarkeit wichtig: Die Zielklasse kann als Begriff wie
\texttt{pipe} angegeben werden, ohne für jede neue lineare Objektart ein
eigenes Modell zu trainieren. Die Videoverarbeitung propagiert die
Objektdefinition über die Sequenz. Damit wird die zeitliche Konsistenz der
Masken zu einem eigenen Qualitätskriterium; eine gute Einzelbildmaske genügt
nicht, wenn das Objekt in späteren Frames verloren geht.

Der kanonische SAM-Ausgang wird als Hierarchie gespeichert. \texttt{default}
ist die unveränderte Objektmaske, \texttt{middle} entsteht durch eine
$5\times5$-Erosion und \texttt{small} durch zwei Erosionsschritte. Die
Erosion wendet ein quadratisches Strukturelement von $5\times5$ Pixeln an:
Ein Pixel gehört nur dann weiterhin zur Maske, wenn die gesamte
$5\times5$-Nachbarschaft maskiert ist; dadurch werden unsichere Randpixel
rund zwei Pixel breit entfernt. Der Wert $5\times5$ ist ein pragmatischer
Standard: Er glättet unscharfe Maskenränder und entfernt dünne
Rauschfragmente, ohne das dünne lineare Objekt über eine sinnvolle Bildbreite
hinaus schrumpfen zu lassen. Die konservativeren Stufen entfernen zusätzlich
unsichere Randpixel, können bei sehr dünnen Objekten aber auch relevante
Objektfläche verlieren. Maskenstufe und optionale Dilatation sind deshalb
explizite Parameter; in der aktuellen Projektfassung wird keine Dilatation
angewendet. Historisch war zwischenzeitlich auch die \texttt{default}-Maske
dilatiert; solche Stände werden als Entwicklungszwischenstände nicht
fortgeführt, sondern durch die verbindliche Festlegung
\texttt{default} = unverändert, \texttt{middle} = einfach erodiert und
\texttt{small} = zweifach erodiert ersetzt (vgl. Abschnitt~\ref{sec:versuchsaufbau}).

\subsection{Structure-from-Motion und Kameramodelle}

Structure-from-Motion schätzt aus korrespondierenden Bildmerkmalen die
Kameraposen und eine sparse 3D-Punktwolke. COLMAP verbindet dafür
SIFT-Merkmalsextraktion, Matching, inkrementelles Mapping und Bundle Adjustment
\parencite{schoenberger2016sfm}. Die Pipeline verwendet unmaskierte Bilder:
Hintergrund- und Ankermerkmale stabilisieren die Orientierung eines glatten,
dünnen Zielobjekts und bleiben für einen späteren GCP-Bezug verfügbar.

Für die Projektarbeit sind drei Kameramodelle relevant. \texttt{PINHOLE}
modelliert getrennte Brennweiten und einen Hauptpunkt, aber keine Verzeichnung. Es setzt damit eine
ideale Bilddomäne voraus. \texttt{SIMPLE\_RADIAL} ergänzt eine gemeinsame
Brennweite um einen radialen Verzeichnungsparameter. \texttt{OPENCV} besitzt
zusätzliche radiale und tangentiale Freiheitsgrade. Das flexiblere Modell ist
nicht grundsätzlich genauer, weil schwache oder ungünstig verteilte Merkmale
zusätzliche Parameter schlecht konditionieren können. Als Produktionsstandard
ist \texttt{OPENCV} gewählt; die Begründung aus der Matrixauswertung steht in
Abschnitt~\ref{sec:produktionskonfiguration}.

Die Modellbezeichnung ist dabei keine hinterlegte Objektivkalibrierung. Das
Skript \texttt{run\_sfm.sh} übergibt COLMAP lediglich die Modellfamilie und mit
\texttt{ImageReader.single\_camera=1} die Annahme eines gemeinsamen
Kameramodells für alle Frames eines Laufs. Eine externe
\texttt{camera\_params}-Datei oder eine Checkerboard-Kalibrierung wird nicht
eingelesen. COLMAP initialisiert die Parameter intern und schätzt Brennweiten,
Hauptpunkt und Verzeichnungskoeffizienten anschließend gemeinsam mit
Kameraposen und 3D-Punkten aus den Bildkorrespondenzen. Die Werte in
\texttt{sparse\_txt/cameras.txt} sind daher lauf- und auflösungsabhängige
Schätzwerte, keine allgemeinen Standarddaten der Kamera.

Bei verzeichnungsbehafteten Modellen bleibt das geschätzte Originalmodell für
SfM und GCP erhalten. Anschließend erzeugt der COLMAP-
\texttt{image\_undistorter} eine ideale PINHOLE-Szene für STS und die
Meshroute. Bei PINHOLE-Eingaben wird am selben Schnittstellenpfad eine
unveränderte Pass-through-Szene bereitgestellt. Diese Unterscheidung ist
entscheidend: Ein direkter PINHOLE-Lauf auf Rohbildern ist eine Ablation der
Modellannahme und kein Nachweis einer erfolgreichen Entzerrung.

\subsection{3D Gaussian Splatting und Segment-then-Splat}

3D Gaussian Splatting stellt eine Szene als Menge anisotroper Gaussians dar
\parencite{kerbl2023gaussian}, das heißt als Ansammlung kleiner, elliptischer
3D-Verteilungen, deren Form und räumliche Ausrichtung pro Gaussian frei
gelernt werden. Ein Gaussian $i$ besitzt unter anderem einen
Mittelpunkt $\boldsymbol{\mu}_i$, eine räumliche Kovarianz
$\boldsymbol{\Sigma}_i$, Opazität $\alpha_i$ und blickabhängige Farbe. Eine
übliche Zerlegung der Kovarianz lautet

\begin{equation}
	\boldsymbol{\Sigma}_i =
	\mathbf{R}_i\,\mathbf{S}_i\mathbf{S}_i^{\top}\mathbf{R}_i^{\top},
	\label{eq:gaussian-covariance}
\end{equation}

wobei $\mathbf{R}_i$ die Orientierung und $\mathbf{S}_i$ die Skalierung
beschreibt. Beim Rendering werden die Gaussians in die Kamera projiziert,
sortiert und über Alpha-Compositing zu einem Bild zusammengeführt. Die
Repräsentation ist deshalb zugleich geometrische Hypothese und renderbare
Radiance-Field-Darstellung.

Segment-then-Splat ergänzt die Gaussians um objektbezogene Zuordnungen
\parencite{lu2025segment}. Das offizielle Verfahren gewinnt seine
Mehrfachansichtsmasken über eine SAM-basierte Tracking-Pipeline
(AutoSeg-SAM2): Ein SAM-Grid-Prompt segmentiert die erste Ansicht, SAM~2
verfolgt das Objekt anschließend über die Sequenz. Nach der Rekonstruktion
wird jedem Objekt ein CLIP-Embedding zugeordnet, über das die Szene offen
vokabularisch durchsucht werden kann. In der Projektpipeline übernehmen die
extern erzeugten SAM-3.1-Masken beide Rollen: Die promptbare temporale
Segmentierung liefert die Objektmasken direkt, und die CLIP-Embedding-
Zuordnung entfällt, weil für den linearen Einobjektfall ausschließlich die
gewählte Objektmaske als Suchbegriff benötigt wird. STS initialisiert die
Szene aus dem COLMAP-Modell und ordnet jeden Gaussian über die vorgegebenen
Masken einem Objekt zu: Während der Objekt-/Maskenphase wird der Bildverlust
auf die maskierte Objektregion der zugeordneten Maske begrenzt, eine
anschließende All-Object-Phase optimiert alle Objekte gemeinsam. In der
Projektpipeline sind diese Masken die bereits segmentierten Eingabebilder in
den drei Erosionsstufen (\texttt{default}, \texttt{middle}, \texttt{small}),
nicht eine vom STS-System erzeugte Nachsegmentierung. Der Projektstandard
umfasst 7000 Gesamtiterationen, davon 5000 innerhalb des Objektcurriculums;
bei einem einzigen Objekt trägt die All-Object-Phase keine zusätzliche
Objektrennung bei, bleibt aber Teil des übernommenen Trainingsschemas.

\subsection{Objektselektion und Meshgewinnung}

Die Ziel-Gaussians werden über die konfigurierte Objekt-ID selektiert; die
Zuordnung stammt aus den STS-Objektdateien (eigene \texttt{obj\_id}-Eigenschaft
der PLY beziehungsweise die Objekt-ID-Tensoren der STS-Ausgabe). Bei einem
Datensatz mit genau einem linearen Objekt ist die Objekt-ID eine
Formalschnittstelle des übernommenen STS-Schemas; die eigentliche Auswahl
trifft der nachgelagerte Opazitäts- und Farbfilter. Ein Opazitätsfilter kann
schwach sichtbare Primitive entfernen. Die in der PLY gespeicherte
Logit-Opazität $o_i$ ist der rohe, unbeschränkte Voraktivierungswert, den das
3DGS-Modell lernt; sie wird erst durch die Sigmoid-Funktion

\begin{equation}
	\alpha_i = \frac{1}{1+\exp(-o_i)}
	\label{eq:opacity}
\end{equation}

in die sichtbare Alpha-Opazität im Bereich $[0,1]$ überführt. Das Filtern über
den Logit erlaubt dabei einen gleichmäßigeren Zugriff auf kleine
Transparenzwerte, als die bereits gesättigte Alpha-Kurve anbieten würde. Nach
Objekt-, Opazitäts- und optionaler Farbauswahl erzeugt die Pipeline eine
getrennte Hochopazitätskopie. Das Setzen auf $\alpha=0{,}999999$ ist eine
Initialisierungspolitik für die geometrische Extraktion, keine physikalische
Aussage über Materialtransparenz: Ein endlicher Wert von exakt 1 ist nach der
Sigmoid-Aktivierung numerisch nicht erreichbar, weil er einen unendlichen
Logit erfordern würde; der nahezu opake Wert stabilisiert die anschließende
Dichte- und Normalenberechnung beim Surface-Sampling.

SuGaR richtet in seiner vollständigen Route Gaussians durch
Oberflächenregularisierung aus und rekonstruiert aus kamerabasierten
Oberflächenstichproben ein Poisson-Mesh\footnote{Die Poisson-
Oberflächenrekonstruktion legt durch eine Punktwolke mit orientierten
Normalen eine glatte Oberfläche, indem sie ein Poisson-Gleichungssystem löst,
dessen Lösungsfeld die gesuchte Indikatorfunktion der Oberfläche beschreibt
\parencite{kazhdan2006poisson}.}
\parencite{guedon2024sugar}. Der lokale Fork kann den
photometrischen Verlust auf die Objektmaske begrenzen. Ein vereinfachter
maskierter RGB-Verlust lautet

\begin{equation}
 \mathcal{L}_{\mathrm{RGB}}^{M} =
 \frac{\sum_p M(p)\,\ell_{\mathrm{RGB}}(p)}
			{\sum_p M(p)+\varepsilon},
 \label{eq:masked-rgb-loss}
\end{equation}

wobei $M$ die ausgewählte Objektmaske (ein binäres Bild) und $p$ eine einzelne
Pixelposition innerhalb dieses Bildes bezeichnen. Der Ausdruck
$\ell_{\mathrm{RGB}}$ kombiniert L1- und DSSIM-Abweichung. Die
Depth-Normal-Consistency der
SuGaR-Coarse-Route vergleicht innerhalb der konservativen Maske aus der
gerenderten Tiefe abgeleitete Normalen mit den Gaussian-Normalen. Im lokalen
Trainingsplan werden diese späten DN-/SDF-Terme erst oberhalb des Zählerstands
9000 aktiviert.

Die Produktionsroute A überspringt die zusätzliche SuGaR-Coarse-Optimierung,
das Gaussian-bound Refinement und das UV-Baking. Sie verwendet die
Original-STS-Gaussians, die projizierte Diamond-Mesh-Tiefe\footnote{Die
Diamond-Mesh-Tiefe wird aus der Projektion der SuGaR-``Diamond''-Primitive
gewonnen: Diese flachen, an den Gaussians ausgerichteten Vierecksflächen
werden in einen Z-Buffer projiziert, aus dem je Pixel eine Tiefe abgelesen
wird; das Surface-Sampling wandelt diese Tiefe anschließend in 3D-Punkte. Die
Alternative ist die Tiefe direkt aus dem Gaussian-Rasterizer.}, Surface-
Sampling, Poisson-Rekonstruktion, Dichtebereinigung, Decimation und optionale
Vertexprojektion. Die Dichtebereinigung entfernt Cluster mit zu wenigen
Stützpunkten, die Decimation reduziert die Dreieckszahl auf das Ziel-Vertice-
Budget, und die optionale Vertexprojektion verschiebt die verbliebenen
Vertices auf die nächstgelegenen Oberflächenstichproben beziehungsweise deren
Tiefenpunkte. Die vollständige SuGaR-Coarse-Route bleibt als
wissenschaftlicher Vergleich erhalten. Damit werden zwei Fragen getrennt:
Wie wird aus einer Gaussian-Wolke ein Mesh gewonnen, und wie verändert eine
zusätzliche Coarse-Optimierung zuvor die Gaussian-Geometrie? SuGaR ist
konzeptionell so ausgelegt, dass Meshes über seine volle Route (Coarse-
Optimierung mit anschließender Oberflächenextraktion) entstehen. Die in der
Projektarbeit verwendete Route A ist deshalb die Abweichung vom
Standardweg: Sie extrahiert direkt aus der unveränderten Original-STS-
Gaussian-Wolke, ohne die Coarse-Optimierung. Diese schlankere Variante wurde
erst im Projektverlauf ergänzend erprobt und hat anschließend die
Produktionsrolle übernommen; die SuGaR-Coarse-Route wird als
Vergleichsarm fortgeführt. Beide Routen bleiben als getrennte und
kombinierbare Bausteine denkbar: Eine spätere Variante könnte auf denselben
Eingaben zusätzlich zur segmentierten Objektrekonstruktion auch eine
vollständige Szene (Gaussians und Mesh) erzeugen, ohne den Objektpfad zu
verändern.

\subsection{Centerline und B-Spline}

DGtal voxeliert das Mesh auf ein diskretes Gitter (Voxelgröße 0,1 in den
lokalen Skaleneinheiten des unskalierten Modells; ohne reale Maßreferenz ist
dies eine relative, keine metrische Größe) und führt eine
topologieerhaltende Skelettierung aus
\parencite{dgtalProject}. Auf verrauschten, dünnwandigen Oberflächen können
dabei Medialflächen (flächige Skelettbereiche, für eine Linientrasse nicht
direkt verwertbar; vgl. die Fußnote in Kapitel~\ref{sec:datengrundlage}) und
zahlreiche kurze Äste entstehen. Der produktive
\texttt{single}-Modus verwendet deshalb den Durchmesserpfad der größten
Skelettkomponente. Dies liefert eine robuste Einzeltrasse, bildet aber keine
verzweigten Netze ab.

Die diskrete Centerline wird durch einen geklemmten uniformen B-Spline des
Grades $p=10$ geglättet. Die Kurve ist

\begin{equation}
	\mathbf{C}(u)=\sum_{i=0}^{n}N_{i,p}(u)\,\mathbf{P}_i,
	\label{eq:bspline}
\end{equation}

mit Kontrollpunkten $\mathbf{P}_i$ und B-Spline-Basisfunktionen $N_{i,p}$.
Der relativ hohe Grad $p=10$ glättet die diskrete Centerline ausreichend,
damit einzelne Ausreißer nicht zu zackigen Kurvensprüngen führen; er ist ein
algorithmischer Glättungsparameter, der im Pipeline-Skript manuell einstellbar
ist, und begründet keine metrische Genauigkeit. Die Anzahl der ausgegebenen
Samples verändert nur die Diskretisierung der Kurve.

\subsection{Objektmaskierte Bildmetriken}

Die Bildmetriken werden objektmaskiert aggregiert. Für $N_M$ gültige
Maskenpixel und normierte RGB-Werte wird der kanalweise gemittelte maskierte
mittlere quadratische Fehler als

\begin{equation}
	\operatorname{MSE}_M=
	\frac{1}{3N_M}\sum_{p:M(p)=1}\lVert I(p)-\hat I(p)\rVert_2^2
	\label{eq:masked-mse}
\end{equation}

berechnet. Daraus folgt bei einem Maximalwert von eins

\begin{equation}
	\operatorname{PSNR}_M=10\log_{10}\!\left(\frac{1}{\operatorname{MSE}_M}\right).
	\label{eq:masked-psnr}
\end{equation}

Ein höherer PSNR-Wert bedeutet eine geringere pixelweise Abweichung. SSIM
bewertet zusätzlich lokale Helligkeit, Kontrast und Struktur; höhere Werte
sind besser \parencite{wang2004ssim}. Im aktuellen Evaluator wird zunächst eine
vollständige SSIM-Karte berechnet und anschließend nur an gültigen
Maskenpixeln aggregiert; der lokale SSIM-Fensterkontext kann daher weiterhin
Randpixel außerhalb der Maske enthalten. LPIPS vergleicht Merkmale eines
vortrainierten Netzes; niedrigere Werte bedeuten eine größere perzeptuelle
Ähnlichkeit \parencite{zhang2018lpips}. Full-frame-Metriken sind
für den object-only Splat deaktiviert, weil sie überwiegend den nicht
rekonstruierten Hintergrund bewerten würden.

PSNR, SSIM und LPIPS sind Ansichtsmetriken. Sie beweisen weder eine korrekte
Meshoberfläche noch eine genaue Centerline. Ebenso garantiert eine vollständige
Maske nicht, dass Kamerageometrie und Surface-Sampling metrisch korrekt sind.
Geometrische RMSE- und Hausdorff-Auswertungen gegen eine unabhängige
GNSS-Referenz bleiben Gegenstand der Bachelorarbeit.
